% Part: applied-modal-logics
% Chapter: temporal-logic
% Section: introduction

\documentclass[../../../include/open-logic-section]{subfiles}

\begin{document}

\olfileid[fr]{aml}{tl}{int}

\olsection{Introduction}

Les logiques temporelles traitent d'énoncés portant sur ce qui sera ou
a été le cas. On attribue à Arthur Prior la création de la logique
temporelle, qu'il appelait \emph{logique des temps verbaux}
(\emph{tense logic}). Notre présentation suivra largement son
approche modale initiale : elle consiste à introduire des opérateurs
temporels dans le cadre de la logique propositionnelle, qui traite en
général les énoncés sans tenir compte de leur temps verbal.

Par exemple, en logique propositionnelle, je pourrais parler d'un
chien, Beezie, qui est tantôt assis, tantôt ne l'est pas, comme il
arrive aux chiens. Il serait contradictoire, en logique classique,
d'affirmer à la fois que Beezie est assis et que Beezie n'est pas
assis. Pourtant, les deux énoncés peuvent être vrais, simplement à
des moments différents ; l'ajout d'opérateurs temporels au langage
permet d'exprimer assez facilement cette situation. Ces opérateurs
permettent aussi de rendre compte de la validité d'inférences telles
que celle qui conduit de «~Beezie recevra une friandise ou une balle~»
à «~Beezie recevra une friandise ou Beezie recevra une balle~».

La logique temporelle soulève toutefois de nombreuses questions
philosophiques qui peuvent nous conduire à préférer un cadre à un
autre. Par exemple, un futur contingent est un énoncé portant sur
l'avenir qui n'est ni nécessaire ni impossible. En disant
«~Richard ira à l'épicerie demain~», nous formulons un énoncé sur
quelque chose qui ne s'est pas encore produit et dont la valeur de
vérité est sujette à discussion. Il est même discutable qu'on puisse
lui \emph{attribuer} une valeur de vérité. Si nous adhérons à un
déterminisme strict, nous pourrions accepter que cet énoncé soit en
fait vrai ou faux, avant même que l'événement en question soit censé
se produire ; il se pourrait simplement que nous ne connaissions
pas encore sa valeur de vérité. Nous pourrions aussi croire à un
avenir véritablement ouvert, dans lequel les valeurs de vérité des
futurs contingents sont indéterminées.

Nos choix de modèles et de cadres de logique temporelle intègrent
nombre de ces conceptions sur la structure et la nature du temps.
Nous pouvons, par exemple, nous demander s'il convient de construire
des modèles dans lesquels le temps est linéaire, ramifié ou même
circulaire. Il peut aussi nous falloir décider si nos modèles
temporels comporteront des instants initiaux et terminaux, et si le
temps sera représenté par des instants discrets ou sous la forme
d'un continuum.

\end{document}
